\documentclass[10pt]{article}
\usepackage[a4paper,margin=1.6cm]{geometry}
\usepackage[T1]{fontenc}
\usepackage{lmodern}
\usepackage{xcolor}
\usepackage{longtable}
\usepackage{array}
\usepackage{tabularx}
\usepackage{colortbl}
\usepackage{booktabs}
\usepackage{ragged2e}
\usepackage{enumitem}
\usepackage{fancyhdr}
\usepackage{hyperref}
\hypersetup{colorlinks=true,urlcolor=blue}
\setlist[itemize]{nosep,leftmargin=*}
\setlength{\parindent}{0pt}
\setlength{\parskip}{4pt}
\pagestyle{fancy}\fancyhf{}\lhead{THE THIRD EYE | Team Parallax}\rhead{14 November 2026 (rev. 3)}\cfoot{\thepage}
\newcommand{\lane}[3]{%
\par\smallskip\noindent\begin{tabularx}{\linewidth}{@{}>{\RaggedRight\arraybackslash}p{2.15cm} >{\RaggedRight\arraybackslash}p{3.65cm} >{\RaggedRight\arraybackslash}X@{}}
\rowcolor{blue!8}\textbf{Agent} & \textbf{Owned area} & \textbf{Assignment} \\
\textbf{#1} & \texttt{\small\detokenize{#2}} & #3 \\
\bottomrule\end{tabularx}\par\smallskip}
\begin{document}
{\LARGE\bfseries THE THIRD EYE}\\[3pt]
{\Large Team missions and AI agent assignments}\\
Five human owners \textbullet{} 15 assigned agent lanes \textbullet{} one evidence system \textbullet{} revision 3

\section*{Purpose and rules}
Build a historical fire calendar that automatically presents qualified comparisons, evidence and limitations. The geographic-access plan is the scientific authority and its \textbf{Weekend Build Brief} sets what is built, in which order, and what is cut first. This document assigns work; it does not certify that results exist or predict an award.

\textbf{Build window (team choice).} All project code, data extracts, fixtures and artifacts are created during the event, November 14--15. Before kickoff the team only does the preparation listed below. This is the team's rule; P5 checks the official legal page and records whether NASA also requires it. Verified NASA rules (11:59 PM local cutoff, teams of at most six at one local event, under-18 media limits) are listed with sources in plan Section 0. Never describe a team choice as a NASA rule.

Each human directs three sessions, each with an assigned task and isolated worktree. Review lanes count; do not force three concurrent writers onto shared files. Assign names, exact paths, base commit, artifacts and test commands before starting. P4 owns shared contracts and serial integration. P1 and P2 approve scientific wording; P5 independently evaluates. Normal permission and event rules apply.

\textbf{Scientific guardrails:} observed fire-cell-days per usable land-cell-days are not counts of unique fires; missing is not zero. Own-support and common-support mappings are distinct. Stability across registered factors is not proof of truth. Geographic subsets require their own eligibility and applicability evidence. Keep sensitivity, applicability and inference status separate. FIRMS hotspot points are a separate view and never enter the grid analysis.

\section*{What changed in this revision}
\begin{itemize}
\item \textbf{Tiers.} Work is ranked P0 (must ship), P1 (ship if P0 is checked) and P2 (only if time). A cut order says what is dropped first when behind.
\item \textbf{Challenge relevance is now P0:} weekly same-sensor calendar, normal season, critical and unusual weeks, selectable sub-areas, and a FIRMS hotspot link with a short explainer.
\item \textbf{Suomi NPP data ceased on November 2, 2026.} The calendar shows the end date. NOAA-20 is the registered continuity path and is a P2 test.
\item \textbf{Static judge site.} Versioned JSON artifacts, no backend, at most 10~MB compressed to start.
\item \textbf{One frozen headline number} with its denominator, plus a plain-language first screen and a first-20-seconds demo.
\item \textbf{Revision 3:} in-app hotspot view moved to P0-late; headline always shows independent units; NASA rules separated from team choices.
\item \textbf{Frontend agent skills} are installed; see \texttt{team-kit/FRONTEND\_SKILLS.md}.
\end{itemize}

\section*{Before kickoff (preparation only, no code or data)}
\begin{itemize}
\item Commit the plan packet and share one agreed commit. Planning documents are not project code.
\item P5: save the official challenge text, resource list and judging guidance under \texttt{challenge/}; read the submission form and video rules.
\item P1/P2: write two or three candidate regions and frozen numeric feasibility rules (seasonal activity in both sensors every training year, a public fire-trend question, clean boundaries, manageable export size). Do not inspect outcomes.
\item Everyone: confirm registration at the same local event, team size (at most six), and any under-18 media rules. P4: deployment and repository access, license choice. P2: Earth Engine and Earthdata sign-in.
\end{itemize}

\section*{Five connected missions}
\renewcommand{\arraystretch}{1.2}
\begin{longtable}{@{}>{\RaggedRight\arraybackslash}p{3cm} >{\RaggedRight\arraybackslash}p{8.8cm} >{\RaggedRight\arraybackslash}p{3cm}@{}}
\rowcolor{blue!12}\textbf{Mission} & \textbf{Outcome and collaboration} & \textbf{Tier} \\
\toprule
M1: Real calendar & P2 extracts NASA data; P1 calculates; P4 serves; P3 displays; P5 checks. Monthly and weekly same-sensor views, normal season, frozen week labels. & P0 \\
M2: Third Eye & P1 calculates Stage A; P2 supplies support; P4/P3 deliver qualified explanations; P5 verifies. & P0 \\
M3: Validation & P5 tests frozen methods against raw, corrected and uncertainty-aware baselines; one frozen headline metric reported with its independent units. & P0 \\
M4: Geography & Primary region plus predeclared raw sub-areas and a FIRMS link at P0; in-app hotspot view P0-late (Sunday 12:00). Extra individually evaluated areas and a wider raw domain are P1. & P0 selector; P1 expansion \\
M5: Jury product & Accessible static site, independently replayed evidence, one-page results, demo script and screen capture. Story is P2. & P0 \\
\bottomrule
\end{longtable}
Missions overlap only where interfaces and dependencies are ready. P0 does not depend on an LLM, IoT, GOES, Stage B or Evidence Repair.

\section*{P1 -- Scientific lead}
\textbf{Accountable for:} the measurement contract, frozen methods and qualified scientific claims; shares approval with P2.
\lane{P1-A}{science/methods/}{Implement own-support $a_H$ and common-support $a_C$ as separate estimands, defined as pooled-rate ratios; Stage A calibration, leave-one-training-year-out variants and policy eligibility.}
\lane{P1-B}{science/calendar/}{Compute monthly and weekly rates, the seasonal baseline, normal season, critical and unusual week rules, period effects and equality boundaries. Apply a practical margin only when justified and registered before outcomes.}
\lane{P1-C}{tests/science/}{Independently check hand-calculated fixtures: missing versus zero, factor invariance, policy compatibility. The composition fixture (8.2/2.8 versus 5.5/5.5) is P1. Review the method without inspecting holdout results.}
\textbf{Also:} with P5, register the block-by-period case list and the headline metric before any holdout is opened. \textbf{Handoff:} versioned method and policy IDs, frozen functions and fixture outputs to P4 and P5. Done when real-input results keep all failed and unavailable states.

\section*{P2 -- NASA data lead}
\textbf{Accountable for:} data access, masks, boundary identity, data quality and source lineage.
\lane{P2-A}{pipeline/masks/ and pipeline/grid/}{Verify product-specific fire and usable classes, UTC alignment, fixed grid, observation support and block-level D/O accounting.}
\lane{P2-B}{pipeline/acquisition/ and pipeline/exports/}{First hour: launch block-level daily exports for the primary and backup regions, both sensors. Split into per-sensor, per-year jobs in this order: 2016--2024 both sensors, then Aqua 2003--2015, then VIIRS 2012--2015. Record job IDs, versions, checksums and failures.}
\lane{P2-C}{data/manifests/}{Source and collection versions, boundaries and sub-area geometry IDs, availability, quality and outage registry (including the Suomi NPP end date), region feasibility, the challenge-to-measurement matrix, the FIRMS link and, from Saturday 20:00, the bounded FIRMS point extract for the in-app hotspot view (max 3 h, credentials server-side, never fed into grid science).}
\textbf{Saturday 13:00:} source access, masks and a real sample are checked; if full history is late, ship the 2016--2024 core and add earlier spans only after each passes its own completeness check. Never show a partial year as complete. \textbf{Handoff:} a verified sample to P1/P4, geometry and status to P3. Write 2020--2021 evaluation outputs to P5-controlled storage; development sessions never read them. P2-C is a source review, not a substitute for independent P5 validation.

\section*{P3 -- Frontend lead}
\textbf{Accountable for:} a clear and accessible application; follow the frontend contract. Start with the \texttt{frontend-design}, \texttt{impeccable}, \texttt{web-design-guidelines}, \texttt{vercel-react-best-practices} and \texttt{playwright-cli} skills. 3D, GSAP, Motion and Blender skills are P2 polish only.
\lane{P3-A}{frontend/calendar/}{Lab shell, area selector with sub-areas, monthly and weekly calendar, normal season, the visible S-NPP end marker and raw/evaluated availability. Own the shared frontend shell and package files.}
\lane{P3-B}{frontend/comparison/}{Third Eye comparison and deterministic summary with separate sensitivity, applicability and inference; accurate failures, units and practical margins. First screen: area, calendar, one finding, one limit and a plain status chip, with no internal schema terms.}
\lane{P3-C}{frontend/evidence/}{Evidence and export views, the FIRMS link with its explainer, the in-app "Reported hotspots" view (P0-late, Sunday 12:00, max 3 h, raw counts only), the landing headline number, and Quick Mission from the same artifact. Keyboard and reduced-motion access. Story is P2.}
\textbf{Handoff:} a geometry, case and version-preserving interface to P4 and P5. Fixtures stay labelled synthetic until replaced by approved real data.

\section*{P4 -- Backend and integration lead}
\textbf{Accountable for:} contracts, CI, deployment, Git integration and stable release artifacts.
\lane{P4-A}{artifact pipeline/}{Versioned static JSON (gzip optional) and explicit unavailable states. A development API is optional; the judge site needs no backend. Measure actual artifact size (target 10~MB compressed) and load time.}
\lane{P4-B}{contracts/}{Freeze Pydantic schemas, frontend types, fixtures, geometry and policy keys and artifact versioning with P1/P2 approval.}
\lane{P4-C}{deploy/ and CI}{Deploy a labelled static skeleton within 60--90 minutes; build checks, public smoke tests and a tested frozen-artifact fallback.}
\textbf{Handoff:} versioned schemas and real-data artifacts to P3; frozen bundles to P5. P4 reviews and integrates approved branches serially, then runs combined checks. Optional InstaBox is a bounded Linux test host only.

\section*{P5 -- Validation and submission lead}
\textbf{Accountable for:} sealed evaluation, independent replay, traceable results and portal submission.
\lane{P5-A}{eval/}{Run the preregistered switch/no-switch experiment over the rule-based block-by-period case list. Report reference-relative errors, answer coverage, abstentions, every registered case and the frozen headline number, always with the independent units (blocks, regions, test years) next to the case count.}
\lane{P5-B}{tests/e2e/}{Check artifact/UI/export parity, geometry and missingness states, and replay the flagship in a clean environment with independent expected results.}
\lane{P5-C}{docs/submission/}{Analyst question and one real published claim, method, source licenses, the one-page results, the demo script and raw screen capture, the claims matrix and the submission form. The user-task rubric is P1.}
\textbf{Holdout rule:} only P5-authorized evaluation sessions can access sealed data in P5-controlled isolation. Preserve the frozen benchmark; any later adaptation gets a new version and never becomes blind retroactively.

\section*{Cross-team handoffs}
\begin{itemize}
\item P2 verified observations $\rightarrow$ P1 calculations $\rightarrow$ P4 artifacts $\rightarrow$ P3 display; P5 independently checks the complete chain.
\item P1/P2 freeze protocol and data membership $\rightarrow$ P5 runs evaluation $\rightarrow$ P4 publishes approved results $\rightarrow$ P3 displays them. P1/P2 review interpretation without retuning the frozen result.
\item P2 geometry and provenance $\rightarrow$ P1 eligibility $\rightarrow$ P4 geometry-specific IDs and cache keys $\rightarrow$ P3 selector $\rightarrow$ P5 cross-region tests.
\item Every handoff includes task ID, base/result commit, plan and contract version, changed paths, actual test results, artifact IDs and limits. Human owners review; P4 integrates.
\end{itemize}

\section*{Tiers and cut order}
\textbf{P0 (all required):} exports launched in the first hour; real-data monthly and weekly calendar with normal season and frozen week labels; primary region plus predeclared raw sub-areas; FIRMS link and explainer; Stage A; static no-backend site; in-app hotspot view (P0-late, Sunday 12:00); frozen held-out report with the headline denominator; independent replay; one-page results, demo script and screen capture; submission.

\textbf{P1 (per owner, after your own P0 is checked):} daily drill-down, additional individually evaluated areas, wider raw domain, expanded diagnostics and the user task. No P1 claim ships without P0 evaluation. Finish and verify by Sunday 18:00.

\textbf{P2 (only if time):} NOAA-20 continuity bridge, recent context, Story, Stage B, Evidence Repair. Each keeps its own gates.

\textbf{Cut order when behind:} P2 first, then P1 (user task and expanded diagnostics, extra areas and wider domain, daily drill-down), then the in-app hotspot view (keep the FIRMS link). Only after that consider other P0 relevance items, and record each unmet challenge item. Never cut a safeguard while still shipping the claim it protects; label that claim unavailable instead.

\section*{Milestones -- local event time}
\begin{longtable}{@{}>{\RaggedRight\arraybackslash}p{3.5cm} >{\RaggedRight\arraybackslash}p{12cm}@{}}
\rowcolor{blue!12}\textbf{Checkpoint} & \textbf{Deliverable and owners} \\
\toprule
\textbf{First hour} & P2 finalizes primary and backup and launches both regions' exports for both sensors. P4 deploys a labelled static skeleton by 90 minutes.\\
\textbf{Saturday 13:00} & Feasibility decision, real sample, and the shorter-span fallback if exports are late: P2 with P1/P4, P3 display.\\
\textbf{Saturday 20:00} & Monthly and weekly calendar, raw area selector, FIRMS link, Stage A and static site integrated: P1--P4.\\
\textbf{Saturday 23:00} & Frozen sensor-switch evaluation, all outcomes, and the headline number on the landing page: P5, with P1/P2 scientific review.\\
\textbf{Sunday 12:00} & Public core, independently replayed evidence bundle (P4/P5) and in-app hotspot view (P2/P3).\\
\textbf{Sunday 18:00} & P1 items verified or cut.\\
\textbf{Sunday 20:00} & One-page results, demo script and raw screen capture recorded: P3/P5.\\
\textbf{Sunday 22:00} & Internal submission and freeze target: whole team, coordinated by P4/P5.\\
\textbf{Sunday 11:59 PM} & Published local submission cutoff; P5 verifies portal timezone and exact cutoff.\\
\end{longtable}
No general 17:00 or 19:00 stop. Checkpoints never authorize fabricated results or weaker scientific gates.

\section*{Short work cycles and conditional scope}
Aim for 5--10 minutes to brief a task, 15--30 minutes for simple setup, 30--60 minutes for a small component and 60--120 minutes for a bounded scientific module to reach first review. These are targets, not validation deadlines. Review small batches frequently; if blocked for 20 minutes, report the dependency and switch to independent work. Data exports and evaluation can take longer.

Optional diagnostics and Evidence Repair share the existing \textbf{six human person-hour} allocation, including integration; start by Sunday 13:00 and retain only checked evidence by 18:00. Recent context and source reproduction use their separate alternative three-hour allocation under the scientific plan. Never present optional features as completed before verification.

\textbf{Release standard:} real data flows end to end, required scientific and UI checks pass, the evidence package can be independently replayed, and all limitations are visible. A polished mockup is not a validated product.
\end{document}
